\clearpage
\begin{summarybox}
\begin{center}\begin{tikzpicture}[font=\small,>=stealth,line join=round]
\draw[gray!40,rounded corners=3pt] (0,0) rectangle (14.2,8.3);\node at (7.1,7.700000000000001) {同过(1,1)，底数跨过1时，上下顺序互换};
\begin{scope}[shift={(1.6625,2.068181818181818)},x=1.8750000000000002cm,y=0.8363636363636363cm]\draw[->,gray] (-0.3,0)--(2.9,0) node[right]{$x$};\draw[->,gray] (0,-0.5)--(0,5) node[above]{$y$};\draw[blue,thick,solid,domain=0.03:2.15,samples=180] plot(\x,{(\x)^2});\draw[teal,thick,dashed,domain=0.03:2.6,samples=180] plot(\x,{\x});\draw[violet,thick,densely dotted,domain=0.03:2.6,samples=180] plot(\x,{sqrt(\x)});\fill[orange!85!black] (1,1) circle(2.2pt);\draw[gray,dashed] (1,0)--(1,4.6);\node[below=5pt] at (1,0){$1$};\end{scope}\node[anchor=west,align=left,text width=5.0cm] at (8.4,4.0) {蓝实线：$x^2$\\[8pt]青虚线：$x$\\[8pt]紫点线：$\sqrt{x}$\\[8pt]$0<x<1$：$x^2<x<\sqrt{x}$\\[8pt]$x>1$：$x^2>x>\sqrt{x}$};
\node[text width=12.6cm,align=center] at (7.1,.65) {这里始终固定系数为1；三条曲线在正半轴上只有公共点(1,1)。};
\end{tikzpicture}\end{center}


  \textbf{\textcolor{CSummary}{一句话核心}}

  幂函数固定系数为 $1$，是为了单独研究指数的作用；指数的正负决定正半轴上的增减方向，而底数与 $1$ 的大小关系决定同底幂的指数比较方向。

  \medskip
  \textbf{\textcolor{CSummary}{快记：一个定点，两个分界}}
  \begin{itemize}[leftmargin=*,itemsep=0.35em]
    \item 定点：所有基本幂函数都过 $(1,1)$。一般式 $Ax^\alpha$ 过 $(1,A)$；基本幂函数在 $(1,1)$ 处的切线斜率就是 $\alpha$。
    \item 指数分界：$0$ 决定上升或下降，$1$ 将正指数的上升走势分成越来越陡、直线、越来越平三种；负指数下降且越来越平，零指数水平。
    \item 底数分界：$x>1$ 时指数越大值越大，$0<x<1$ 时指数越大值越小；$x=1$ 时都等于 $1$。
  \end{itemize}

  \medskip
  \textbf{\textcolor{CSummary}{快懂：为什么？}}

  $1$ 是乘法中不改变数值的数。先固定系数，才能只比较指数；先研究 $x^2$、$x$、$1/x$ 等基本图像，就能理解其他由它们伸缩或翻转得到的图像。指数比较的实质是多乘了一个大于 $1$ 或小于 $1$ 的正数；负指数的实质是取倒数。

  \medskip
\end{summarybox}
\clearpage
\begin{summarybox}
  \textbf{\textcolor{CSummary}{快用：解题顺序}}
  \begin{enumerate}[leftmargin=*,itemsep=0.3em]
    \item 题目说是幂函数：把 $x^\alpha$ 前面的系数化为 $1$，有参数就先求参数。
    \item 判断定义域，确认底数能否为零或负数。
    \item 同指数比较底数：看指数正负；同底数比较指数：看底数与 $1$ 的大小关系。
    \item 遇到负指数可转成倒数；两边分居 $1$ 的两侧时，可用 $1$ 作比较的桥梁。
    \item 解不等式最后检查定义域，排除无意义的点。
  \end{enumerate}

  \medskip
  \textbf{\textcolor{CSummary}{使用边界}}

  五种走势与同底实数幂比较的统一规则只针对 $x>0$。其他部分必须根据具体指数的定义域与奇偶性另行判断。$\alpha=0$ 的水平线在 $(0,1)$ 处有空缺；只有正指数的完整幂函数图像经过原点。
\end{summarybox}
